\section{A product rule for the plaque probabilities}\label{sec:product}

The projective measures constructed in the preceding section retain
information obtained by conditioning the measures \(\mu_i\) at depths
tending to infinity.  We now prove a product rule for these measures.
The factors in the rule belong to the \emph{same augmented state}; this
point will allow us to compose root directions in the next section.

Product decompositions for genuine leafwise conditional measures were
developed by Einsiedler and Katok
\citep[Proposition~8.3]{EK2003}
\citep[Theorems~8.4--8.5]{EK2005}; see also
Einsiedler and Lindenstrauss
\citep[Corollary~8.8 and Theorem~9.8]{EL2010}.
The use of a diagonal element that contracts one factor and fixes
another follows this established method.  Here the measures are the
augmented-state data obtained by conditioning before weak convergence.
We prove the product identity for those data using the preceding
mesh comparisons; the cited leafwise theorems are not applied to them.

We continue to use the notation \(p\prec q\) for \(p_i/q_i\to0\).
Every scale used below lies in the plateau
\([l_{-,i},l_{+,i}]\), with room for the displayed separations and for
additional separations at both ends.  Expectations in a chart mean
integration against the restriction of \(\mu_i\) to that chart; thus
all error estimates are weighted by the mass of the chart.  A chart of
zero limiting mass can be omitted.

\begin{theorem}[Product rule]\label{thm:product}
Let \(V\) and \(W\) be pattern groups, let \(U=U_\alpha\) be a
one-dimensional root group, and suppose that multiplication gives a
coordinate decomposition
\[
                 V=WU,\qquad W\mathrel{\triangleleft}V.
\]
Suppose also that some \(a_0\in A\) fixes \(U\) by conjugation and that
the automorphism
\[
                 \Delta(w)=a_0^{-1}wa_0
\]
strictly contracts \(W\).  Then, for almost every augmented state
\(\omega\),
\begin{equation}\label{eq:product-rule}
 m^V_\omega
   \ \propto\
   ( (w,u)\longmapsto wu)_*
          \bigl(m^W_\omega\otimes m^U_\omega\bigr).
\end{equation}
Here \(\propto\) denotes equality up to multiplication by a positive
constant.
\end{theorem}

The case \(W=\{1\}\) follows from the plaque comparisons.  Henceforth
\(W\ne\{1\}\).  The root constraints defining a contracting logarithm
are rational linear equalities and strict inequalities.  We can
therefore choose \(a_0\), and rescale its logarithm, so that in the
matrix-entry coordinates of \(W\),
\begin{equation}\label{eq:integer-contraction}
       \Delta(w)_{jk}=2^{-d_{jk}}w_{jk},
       \qquad d_{jk}\in\N,\quad d_{jk}>0.
\end{equation}
Whenever \(jk\) and \(k\ell\) compose, their weights add:
\(d_{j\ell}=d_{jk}+d_{k\ell}\).
This is also the compatibility that makes \(\Delta\) an automorphism.

We use the simultaneous charts in the stock of the preceding section:
\begin{equation}\label{eq:simultaneous-product-chart}
        (q,w,u)\longmapsto\psi(q)wu
\end{equation}
on product domains.  The transverse coordinates are \(q\) for \(V\),
\((q,w)\) for \(U\), and \((q,u)\) for \(W\).  The last assertion uses
normality of \(W\); its plaque parametrization may be a smooth
reparametrization of the group coordinate.  Compatible transverse
grids are used when the same \(q\) occurs in two conditionings.

The proof has two parts.  Contraction of \(W\) shows that the \(U\)
coordinate has the same law in the \(V\) and \(U\) data, independently
of the \(W\) coordinate.  Identifying the remaining factor requires a
different argument.  A depth-\(S\) \(U\)-cell label contains only
\(O(S)\) information.  When \(S\prec T\), revealing this label can
appreciably change the law of a fixed-length further \(W\)-tile
refinement, conditional on its parent tile, at only a vanishing
proportion of \(T\) successive stages.  Expanding a selected parent
tile then turns this refinement comparison into a comparison of
bounded \(W\) tests.

\subsection{The factor fixed by the contraction}

\begin{lemma}\label{lem:u-factor}
In each chart \eqref{eq:simultaneous-product-chart}, the limiting
\(V\)-plaque probability is a product in the coordinates \(w,u\).
Its \(u\) marginal agrees with the \(U\)-plaque probability at the
sampled state.  Consequently there is a projective locally finite
measure \(\nu^W_\omega\) on \(W\) such that
\begin{equation}\label{eq:provisional-product}
 m^V_\omega
 \ \propto\
 ((w,u)\longmapsto wu)_*
          (\nu^W_\omega\otimes m^U_\omega).
\end{equation}
\end{lemma}

\begin{proof}
Choose interior scales \(l_-\prec T\prec F\prec l_+\), and apply
the diagonal motion \(a_0^{\lfloor T\rfloor}\).
In one fixed input chart, a patch specifies a \(q\)-cell at depth
\(F\), with the whole \(w\) and \(u\) ranges.  A subpatch also
specifies a \(w\)-cell at depth \(F\).  Denote these two labels by
\(Q_F\) and \(W_F\).

Fix an arbitrarily small permitted coverage loss.  Take finitely
many output \(U\)-charts, with inner boxes and sufficiently large
plaque extensions, covering all but that loss of output probability.
The law of the output point is still \(\mu_i\).  An input patch whose
image meets one of these inner boxes lies in the corresponding
extended chart, for all sufficiently large \(i\): its \(U\) range
has bounded fixed size, its \(W\) displacements shrink by
\(O(e^{-cT})\), and its \(q\) variation contributes at most
\[
                       O(2^{-F}e^{CT})
\]
for fixed constants \(c,C>0\).  Assign such patches to output charts
and omit the others.  The omitted mass is bounded by the coverage
loss.

In an assigned \(U\)-chart, the transverse image of a patch has
diameter
\[
                       O(e^{-cT}+2^{-F}e^{CT})
                       =o(2^{-l_-}).
\]
The transverse projection is independent of \(u\).  Its restriction
to the input \((q,w)\) range is injective: the chart has a common
open \(U\) range for all these inputs, and two equal transverse
coordinates would otherwise give overlapping images of two
different points of the injective input chart.  In the output
plaque coordinate, the input \(u\) interval is translated by a
parameter varying by \(o(1)\) across the patch.  Thus patches and
their subpatches admit a common collection of inner \(U\) bins.

Here are the trimming details needed for
Lemma~\ref{lem:patch-comparison}.  A patch meeting a coarse output
wall has its sampled point in a collar of relative width \(o(1)\);
the grid convention makes the total mass of these patches \(o(1)\).
Fine output transverse cells can be pulled back at cost at most
\(e^{CT}\).  Since \(l_+\succ F\), the pullback variation is
negligible compared with the input \(q\)- and \(w\)-meshes.
Away from their grid walls, the whole fine cell therefore stays
inside the prescribed input labels.  Fixed small collars of the
input chart boundary have arbitrarily small mass.  After removing
these collars, sufficiently fine fixed output \(U\) bins lie wholly
in the admitted plaque interval, except for arbitrarily small
sampling mass.  One may first trim each subpatch into products of
fine transverse cells and these common plaque bins, and take their
union as a trim of the patch.  Both levels of the comparison use the
same coarse output cell and the same plaque bins.

The patch/subpatch conclusion of
Lemma~\ref{lem:patch-comparison} now gives, for every continuous
function \(f\) on the closure of the input \(u\) interval,
\begin{equation}\label{eq:u-independence-prelimit}
 \int
 \left|
   \E_{\mu_i}\bigl(f(u)\mid Q_F,W_F\bigr)
    -\E_{\mu_i}\bigl(f(u)\mid Q_F\bigr)
 \right|\,\dd\mu_i
                         \longrightarrow0.
\end{equation}
To obtain the limit zero, first make the bin and chart-collar losses
arbitrarily small for the fixed finite cover, and then make the
coverage loss arbitrarily small.  The common translation of the
\(u\) coordinate causes no extra error, because its variation tends
to zero and \(f\) is uniformly continuous.

The plateau comparison replaces the displayed depths by the depths
defining the state data.  It follows that the \(u\) marginal of the
\(V\) data is the \(U\) probability at the sample.  For independence,
let \(g\) be a bounded continuous function of \(w\), with modulus
of continuity \(\omega_g\).  Approximate it by a constant
function \(g_F\) on each \(W_F\)-cell.  Conditional expectation,
followed by the triangle inequality, gives
\[
\begin{split}
 &\int\left|
   \E_{\mu_i}(f(u)g(w)\mid Q_F)
     -\E_{\mu_i}(f(u)\mid Q_F)\E_{\mu_i}(g(w)\mid Q_F)
                  \right|\,\dd\mu_i\\
 &\quad\le
 \|g\|_\infty
 \int\left|
   \E_{\mu_i}(f(u)\mid Q_F,W_F)
       -\E_{\mu_i}(f(u)\mid Q_F)
       \right|\,\dd\mu_i
       +2\|f\|_\infty\omega_g(C2^{-F}).
\end{split}
\]
The right side tends to zero by
\eqref{eq:u-independence-prelimit}.  Passing to the state limit
therefore gives
\[
 P^V(f(u)g(w))=P^V(f(u))P^V(g(w))
\]
where \(P^V\) denotes the limiting \(V\)-chart probability.
A countable dense family of continuous
tests proves that the limiting probability is a product.

For a sampled plaque point with coordinates \(w_s,u_s\), the group
displacement from that point to \(wu\) has \(W,U\) coordinates
\begin{equation}\label{eq:product-displacement}
 \left(\Ad_{u_s^{-1}}(w_s^{-1}w),\ u_s^{-1}u\right).
\end{equation}
This changes the two coordinates separately, so the product
assertion holds in displacement coordinates as well.
Proposition~\ref{prop:plaque-comparison} identifies overlapping
chart probabilities projectively.  Apply this comparison on
product ranges containing the identity and exhausting \(W\times
U\); charts with the required plaque extensions are available.
The resulting \(W\) marginals agree projectively and define
\(\nu^W_\omega\).  This proves \eqref{eq:provisional-product}.
\end{proof}

It remains to prove that \(\nu^W_\omega=m^W_\omega\) projectively.
The contraction just used does not shrink the \(U\) direction, so
interchanging the names of the two factors would not prove this.
We instead compare small \(W\) pieces before a reverse diagonal
motion with bounded \(W\) ranges after that motion.

\subsection{Nested tiles and their information cost}

We first construct subdivisions of \(W\) compatible with \(\Delta\).
Let \(\Gamma_W\) be the lattice of pattern matrices with integral
entries.  Haar measure on \(W\) is Lebesgue measure in its matrix
entries, with \(\Gamma_W\) of covolume one.

\begin{lemma}[Tiles adapted to the contraction]\label{lem:tiles}
There is a compact set \(\mathcal F\subset W\), of Haar measure one
and with Haar-null boundary, such that the translates
\(\gamma\mathcal F\), \(\gamma\in\Gamma_W\), tile \(W\) with disjoint
interiors.  Moreover,
\[
 \mathcal F
    =\bigcup_{d\in\mathcal D}\Delta(d)\Delta(\mathcal F),
 \qquad
 \mathcal D=
 \{\gamma\in\Gamma_W:0\le \gamma_{jk}<2^{d_{jk}}
                         \text{ for each pattern entry }jk\},
\]
where \(d_{jk}\) is the contraction weight in
\eqref{eq:integer-contraction}.
The union is disjoint off boundaries.

For every integer \(k\ge0\), the tiles
\[
       \gamma\Delta^k\mathcal F,\qquad
       \gamma\in\Delta^k\Gamma_W,
\]
form a subdivision, and these subdivisions are nested as \(k\)
increases.  They may all be shifted by the same left translation
\(g\).  If \(g\) is distributed according to Haar probability on
\(W/\Gamma_W\), then, for every fixed \(w\in W\), its coordinate
in its shifted level-\(k\) tile has uniform Haar distribution on
\(\mathcal F\), after rescaling by \(\Delta^{-k}\).
In particular, the probability of lying in a fixed
\(\rho\)-collar of \(\partial\mathcal F\) in these coordinates
tends to zero as \(\rho\downarrow0\), uniformly in \(k\) and \(w\).
\end{lemma}

\begin{proof}
For this proof write \(d_\alpha\) for the positive integral weight
of a pattern entry \(\alpha\), and put \(B_\alpha=2^{d_\alpha}\).
The digit set consists of integral pattern matrices \(d\) whose
\(\alpha\) entry belongs to \(\{0,\ldots,B_\alpha-1\}\).
Define
\begin{equation}\label{eq:tile-digits}
 \mathcal F=
 \left\{
      \prod_{\nu=1}^{\infty}\Delta^\nu(d_\nu):
                      d_\nu\in\mathcal D
 \right\},
\end{equation}
with factors ordered by increasing \(\nu\).
Contraction and the finite nilpotent multiplication formulas imply
uniform convergence of these products.  More explicitly, in any
entry the product is a finite sum of convergent series of products
of digit entries, with positive geometrically decaying weights.
Thus \(\mathcal F\) is compact.

Order the pattern entries by increasing root height after putting
the pattern in upper triangular form.  In an entry \(\alpha\), the
coordinate in \eqref{eq:tile-digits} is
\[
       \sum_{\nu\ge1}B_\alpha^{-\nu}(d_\nu)_\alpha
                  +C_\alpha,
\]
where \(C_\alpha\) is determined entirely by digits in entries of
smaller height.  At height one, this is the usual base-\(B_\alpha\)
expansion of a number in \([0,1]\).  Inductively, once the
lower-height digits have been specified outside their null sets
of ambiguous expansions, the possible \(\alpha\) coordinates
form the interval \([C_\alpha,C_\alpha+1]\), with a unique digit
expansion except at its countably many base endpoints.

Left multiplication by an integral pattern matrix has the same
triangular form: the \(\alpha\) coordinate changes by its integral
\(\alpha\) entry plus a function of already chosen lower-height
entries.  Reduce entries in this order.  At each step exactly one
integer places the remaining coordinate in an interval of length
one, except at an endpoint, and its base expansion is then unique
almost everywhere.  Fubini's theorem at each of the finitely many
entries proves that almost every point of \(W\) belongs to exactly
one translate \(\gamma\mathcal F\).  It also proves that
\(\mathcal F\) has Haar measure one.

The translates are locally finite, because \(\mathcal F\) is
compact and \(\Gamma_W\) is discrete.  Their union is consequently
closed, and the almost-everywhere covering just proved implies
that they cover all of \(W\).  Splitting the first digit from
\eqref{eq:tile-digits} gives
\begin{equation}\label{eq:tile-subdivision}
          \mathcal F
             =\bigcup_{d\in\mathcal D}
                    \Delta(d)\Delta(\mathcal F).
\end{equation}
Triangular reduction also shows that \(\mathcal D\) represents the
right cosets of \(\Delta^{-1}\Gamma_W\) in \(\Gamma_W\):
\[
      \Gamma_W=
          \bigsqcup_{d\in\mathcal D}
                       (\Delta^{-1}\Gamma_W)d.
\]
Applying \(\Delta\) shows that the origins in the subdivision of
all level-zero tiles are exactly
\[
            \Delta\Gamma_W
                  =\bigsqcup_{d\in\mathcal D}\Gamma_W\Delta(d).
\]
Iteration proves the asserted level-\(k\) subdivision and nesting.

For completeness, these compact tiles have the regularity needed
for boundary trimming.  The closed translates cover \(W\), so the
Baire category theorem gives nonempty interior to \(\mathcal F\).
Every point of \(\mathcal F\) lies in a descendant tile in the
iterated subdivision \eqref{eq:tile-subdivision}.  Such descendants
have nonempty interior and diameters tending to zero.  Therefore
\(\mathcal F\) is the closure of its interior.  An interior
intersection of two translates would then give a positive-volume
overlap, contradicting almost-everywhere uniqueness.  A boundary
point of one translate belongs to another: otherwise local
finiteness would give a neighborhood covered only by the first
translate, making that point interior.  Thus the boundary is
contained in the union of the pairwise overlaps, which is null.
Boundary assignments can be chosen measurably and consistently
through the nested subdivisions.

It remains to check the law of a shifted tile coordinate; this
check matters because the sampling measure on \(W\) need not have
a density.  Since the weights are positive integers,
\(\Gamma_W\subset\Delta^k\Gamma_W\).
The projection
\[
                 W/\Gamma_W\longrightarrow W/\Delta^k\Gamma_W
\]
sends Haar probability to Haar probability.  Inversion and right
multiplication by \(w\) send the latter distribution to invariant
volume on the left quotient
\(\Delta^k\Gamma_W\backslash W\).
The representative of \(g^{-1}w\) in \(\Delta^k\mathcal F\)
therefore has normalized Haar distribution on that tile, and
rescaling gives Haar measure on \(\mathcal F\).
This is precisely the coordinate of \(w\) in its tile
\(g\gamma\Delta^k\mathcal F\).
The null boundary implies the collar assertion by continuity from
above.  Fubini's theorem permits the same conclusion after
sampling \(w\) from any probability measure.
\end{proof}

For finite labels we write \(H\) for Shannon entropy, with natural
logarithms, and
\[
       I(Y;Z\mid X)=H(Y\mid X)-H(Y\mid X,Z)
\]
for conditional mutual information.  Tile labels restricted to a
fixed chart are finite at every fixed level.

\begin{lemma}[Information in a transverse label]\label{lem:stage-information}
Fix a chart \eqref{eq:simultaneous-product-chart} and integers
\(F,S\ge1\).  Let \(Q_F\) be its depth-\(F\) \(q\)-cell label,
and let \(U_S\) be its depth-\(S\) \(u\)-cell label.
Let \(K_k\) be the level-\(k\) tile label of \(w\), with any common
shift as in Lemma~\ref{lem:tiles}, restricted to the chart.
For every pair of integers \(T\ge j\ge1\),
\begin{equation}\label{eq:stage-information}
 \sum_{k=0}^{T-j}
       I_{\mu_i}(K_{k+j};U_S\mid Q_F,K_k)
                           \le Cj(1+S).
\end{equation}
Here the measure in the chart is normalized for defining entropy;
its mass is restored when errors are integrated over the chart.
The constant \(C\) depends only on the fixed coordinate domain,
and is independent of \(i,F,S,T,j\), the sampling measure, and
the tile shift.
Consequently, along any sequence with \(T\to\infty\) and
\(S/T\to0\), for every fixed \(j\), averaged over
\(T/2\le k\le3T/4\), revealing \(U_S\) changes the conditional
law of \(K_{k+j}\), given \((Q_F,K_k)\), by \(o(1)\) in expected
\(\ell^1\) distance.  The same assertions hold after averaging
over shifts.
\end{lemma}

\begin{proof}
The number of \(u\)-cells meeting its fixed bounded interval is
\(O(2^S)\), so \(H(U_S)\le C(1+S)\).
Because the tile labels are nested, the chain rule gives
\[
 \sum_{k=0}^{L-1}
       I(K_{k+1};U_S\mid Q_F,K_k)
       = I(K_L;U_S\mid Q_F,K_0)
       \le H(U_S).
\]
For a jump of length \(j\), expand its mutual information into
these one-stage increments.  Each increment occurs at most \(j\)
times in the sum over starting stages, proving
\eqref{eq:stage-information}.  Pinsker's inequality and
Cauchy--Schwarz bound the average conditional \(\ell^1\) error.
For \(T\ge4j\), write
\(\mathcal K_T=\{k\in\Z:T/2\le k\le3T/4\}\).  The bound is
\[
 \begin{split}
 &\frac1{|\mathcal K_T|}\sum_{k\in\mathcal K_T}
 \E\bigl\|
   \mathcal L(K_{k+j}\mid Q_F,K_k,U_S)
       -\mathcal L(K_{k+j}\mid Q_F,K_k)
 \bigr\|_1\\
 &\hspace{25mm}\le C\sqrt{\frac{j(1+S)}{T}}.
 \end{split}
\]
This tends to zero for fixed \(j\).  The entropy bound is uniform
in the shift, so it can also be averaged over shifts.
\end{proof}

\subsection{Identification of the contracted factor}

We now compare \(\nu^W_\omega\) with \(m^W_\omega\).  Choose
interior scales
\begin{equation}\label{eq:w-factor-scales}
                       l_-\prec S\prec T\prec F\prec l_+.
\end{equation}
Fix nonnegative functions \(f,\phi\in C_c(W)\), bounded by one,
with \(\phi(1)>0\).  It suffices to show
\begin{equation}\label{eq:w-ratio-target}
        \frac{\nu^W_\omega(f)}{\nu^W_\omega(\phi)}
                  =
        \frac{m^W_\omega(f)}{m^W_\omega(\phi)}
\end{equation}
almost surely.  Identity support makes the denominators positive.
A countable determining family of such tests will then give
projective equality.

The comparison is made at the output of the motion
\begin{equation}\label{eq:reverse-motion}
                         a_0^{-(k+m_0)},
\end{equation}
where \(T/2\le k\le3T/4\) and \(m_0\) is a fixed positive integer.
A level-\(k\) input tile has, relative to its translated origin,
the fixed output shape \(\Delta^{-m_0}\mathcal F\).
The role of \(m_0\) is to make this shape contain the supports of
the tests about most sampled points.  The role of the much finer
level \(k+j\), with \(j\) fixed after \(m_0\), is to resolve those
tests at output.

The two distributions to be compared can now be specified before
the geometric details.  Inside one input chart, a parent patch has
label \(p=(Q_F,K_k)\), leaving the whole \(u\) interval available;
a subpatch has label \(p'=(Q_F,K_k,U_S)\).
Lemma~\ref{lem:stage-information} compares
\begin{equation}\label{eq:w-two-tile-laws}
 \mathcal L_{\mu_i}(K_{k+j}\mid Q_F,K_k)
 \quad\text{and}\quad
 \mathcal L_{\mu_i}(K_{k+j}\mid Q_F,K_k,U_S).
\end{equation}
The measure here is restricted to the input chart.
After \eqref{eq:reverse-motion}, the parent \(W\) tile has fixed
shape \(\Delta^{-m_0}\mathcal F\), and its descendants have shape
\(\Delta^{j-m_0}\mathcal F\).  Increasing the fixed \(j\) therefore
allows \eqref{eq:w-two-tile-laws} to compare bounded continuous
\(W\) tests on that output range.

In a simultaneous output chart with coordinates \((q',w',u')\),
the transverse coordinate of the whole patch is \(q'\);
for its thin \(u\)-subpatch it is \((q',u')\).
The geometry and plateau estimate below will consequently compare
the whole-patch test law with a restriction of the \(V\) data,
and the subpatch test law with a restriction of the \(W\) data.
These comparisons have normalization factors and retain an inner
\(u'\)-range in the \(V\) data.  We will first compare their cross
products, then use Lemma~\ref{lem:u-factor} to remove that \(u'\)
restriction, and finally control the denominators.  This is how
the two conditional tile laws will identify the \(W\) factor
without presupposing the desired projective equality.

Figure~\ref{fig:product-mechanism} contrasts these two comparisons:
contraction determines the \(U\) factor, whereas expansion must be
combined with the information estimate to identify the \(W\) factor.
\begin{figure}[tbp]
\centering
\begin{tikzpicture}[
  x=1cm,y=1cm,
  font=\small,
  >={Stealth[length=2mm]},
  patch/.style={draw=black!75,line width=.55pt,fill=black!4},
  selected/.style={draw=blue!65!black,line width=.55pt,fill=blue!16},
  coarse/.style={draw=black!70,dashed,line width=.65pt},
  fine/.style={draw=black!35,line width=.25pt},
  bins/.style={draw=black!70,densely dotted,line width=.5pt},
  note/.style={align=left,font=\footnotesize},
  leader/.style={draw=black!65,line width=.4pt,->}
]
\path[use as bounding box] (0,-5.15) rectangle (14.2,4.35);

\node[anchor=west,font=\small\bfseries] at (0,4.1)
  {(a) Contract $W$: identify the $U$ factor and prove independence};
\node[align=center,font=\footnotesize] at (2.35,3.67)
  {one $q$-cell; full $w,u$ ranges};
\draw[patch] (.65,1.3) rectangle (4.05,3.35);
\draw[selected] (2.15,1.3) rectangle (2.75,3.35);
\node[font=\footnotesize,anchor=north] at (2.45,1.19) {$w$-cell};
\node[anchor=east] at (.51,3.2) {$u$};
\node[anchor=north] at (3.9,1.19) {$w$};

\draw[->,line width=.8pt] (4.65,2.35) -- (7.1,2.35);
\node[above] at (5.87,2.45) {$a_0^{t}$, $t\asymp T$};
\node[below,font=\footnotesize,align=center] at (5.87,2.18)
  {$W$ contracts; $U$ is fixed\\$T\prec F$};

\draw[coarse] (7.95,1.15) rectangle (10.15,3.5);
\draw[patch] (8.55,1.3) rectangle (9.55,3.35);
\draw[selected] (8.991,1.3) rectangle (9.168,3.35);
\foreach \x in {8.6,8.7,8.8,8.9,9.0,9.1,9.2,9.3,9.4,9.5}
  \draw[fine] (\x,1.3) -- (\x,3.35);
\foreach \y in {1.78,2.31,2.84}
  \draw[bins] (7.95,\y) -- (10.15,\y);
\node[font=\footnotesize] at (9.05,3.73) {output $U$ chart};
\node[note,anchor=west] at (10.58,3.18)
  {One coarse transverse cell\\contains the whole patch.};
\draw[leader] (10.5,3.15) -- (10.17,3.15);
\node[note,anchor=west] at (10.58,2.26)
  {Fine transverse cells\\resolve the $w$-label.};
\draw[leader] (10.5,2.24) -- (9.28,2.24);
\node[note,anchor=west] at (10.58,1.42) {Common $U$ bins};
\draw[leader] (10.5,1.48) -- (10.04,1.78);
\node[align=center,font=\footnotesize] at (7.1,.5)
  {Mesh comparison: revealing the $w$-cell changes the $U$-bin law by $o(1)$ in sampling mean.};

\node[anchor=west,font=\small\bfseries] at (0,-.22)
  {(b) Expand a $W$ tile: identify the remaining $W$ factor};
\node[align=center,font=\footnotesize] at (2.35,-.86)
  {one stage-$k$ tile\\shape $\Delta^k\mathcal F$};
\draw[patch] (2.0,-3.7) rectangle (2.7,-1.55);
\draw[selected] (2.0,-2.68) rectangle (2.7,-2.43);
\foreach \x in {2.175,2.35,2.525}
  \draw[bins] (\x,-3.7) -- (\x,-1.55);
\node[anchor=east] at (1.86,-1.72) {$u$};
\node[anchor=north] at (2.65,-3.81) {$w$};
\node[note,anchor=east] at (1.72,-2.57) {$u$-cell\\depth $S$};
\draw[leader] (1.8,-2.57) -- (1.98,-2.57);

\draw[->,line width=.8pt] (3.65,-2.57) -- (6.5,-2.57);
\node[above] at (5.08,-2.45) {$a_0^{-(k+m_0)}$};
\node[below,font=\footnotesize,align=center] at (5.08,-2.73)
  {$W$ expands; $U$ is fixed\\$S\prec T\prec F$, $k\asymp T$};

\node[align=center,font=\footnotesize] at (8.63,-.88)
  {fixed output shape\\$\Delta^{-m_0}\mathcal F$};
\draw[patch] (7.05,-3.7) rectangle (10.2,-1.55);
\draw[selected] (7.05,-2.68) rectangle (10.2,-2.43);
\foreach \x in {7.8375,8.625,9.4125}
  \draw[bins] (\x,-3.7) -- (\x,-1.55);
\draw[coarse] (6.9,-2.92) rectangle (10.35,-2.2);
\foreach \y in {-2.63,-2.58,-2.53,-2.48}
  \draw[fine] (7.05,\y) -- (10.2,\y);
\node[note,anchor=west] at (10.58,-1.94)
  {Only the $u$-subpatch fits\\in one coarse transverse\\cell for the $W$ chart.};
\draw[leader] (10.5,-2.02) -- (10.35,-2.2);
\node[note,anchor=west] at (10.58,-3.18)
  {Fine transverse cells\\resolve the $u$-label.};
\draw[leader] (10.5,-3.05) -- (9.72,-2.6);
\node[align=center,font=\footnotesize] at (8.63,-4.04)
  {Stage-$(k+j)$ subtiles resolve $W$ tests.};
\node[align=center,font=\footnotesize,text width=14cm] at (7.1,-4.78)
  {Information comparison: an $O(S)$ budget over $T$ stages makes the $u$-label change\\
   the descendant-tile law by $o(1)$ in sampling mean, at a suitable stage $k\asymp T$.};
\end{tikzpicture}
\caption{The two different comparisons in the product rule. Each input picture
fixes a $q$-cell at depth $F$; the coordinate $q$ is suppressed. Dashed outlines
denote coarse output transverse cells; thin solid lines denote fine transverse
cells. Dotted lines denote fixed $U$ bins in (a) and stage-$(k+j)$ $W$ subtiles in (b).
Shading marks a labelled subpatch, not a density. In (a), contraction permits
one coarse transverse-cell comparison for the whole patch and its subpatches.
In (b), expansion leaves the full $U$ interval broad. The full patch is compared
with $V$ data and the thin $u$-subpatch with $W$ data; the information estimate
compares their descendant-tile laws and hence their resolved $W$ tests. All widths and
tile shapes are schematic: $W$ can have several dimensions and its tiles need
not be rectangles. The resulting factor identifications are projective.}
\label{fig:product-mechanism}
\end{figure}


\paragraph{Coverage and the order of choices.}
Fix a permitted exceptional probability \(\epsilon>0\).
Take finitely many input charts
\eqref{eq:simultaneous-product-chart}, with inner boxes and fixed
margins, covering all but a sufficiently small multiple of
\(\epsilon\).  Estimates can be added over these charts; their
finite multiplicity only changes constants.
Use the patches and subpatches just described, with whole
level-\(k\) \(w\)-tiles.  Tiles meeting
the fixed inner boxes lie inside the outer input \(w\) ranges for
large \(i\), because their diameters tend to zero.

First choose a tile-coordinate collar thin enough that its
expected mass is much smaller than the permitted coverage loss.
Lemma~\ref{lem:tiles} does this uniformly in \(k\) and in the input
distribution.  Outside this collar the sample has a fixed
positive margin in \(\mathcal F\).  Choose \(m_0\) so large that
every displacement in a fixed compact set containing
\(\supp f\cup\supp\phi\) and the identity lies, with margin, in
the output tile about such a sample.
Indeed, pull such a displacement back to the input tile
coordinate.  It is contracted by \(\Delta^{m_0}\); the
conjugation coming from the input \(u\) coordinate is uniformly
bounded, and commutes with \(\Delta\), since \(a_0\) fixes \(U\).
Thus the pulled-back displacement fits inside the chosen margin
for sufficiently large fixed \(m_0\).
The input inner \(u\) intervals also give a fixed \(U\)
neighborhood about the sampled output point.

Now take finitely many output charts of the same simultaneous
form, with inner boxes and \(V\)-extensions large enough for the
images of these patches.  Such a finite cover loses an arbitrarily
small further probability.  For fixed \(m_0\), all displacements
from a sample to points in its output patch lie in a fixed compact
range: use the tile origin as base in \(W\), a fixed base in \(U\),
and the fixed shape \(\Delta^{-m_0}\mathcal F\).
The \(q\) variation after \eqref{eq:reverse-motion} is
\(O(2^{-F}e^{CT})\).
Hence a patch meeting an output inner box fits in its extended
chart with margin.  We can assign these patches to the finitely
many output charts, or perform the estimates for each applicable
input--output pair.  We keep only samples with the indicated
inner margins.

These choices fix the chart cover, the output tile extent, and
the test-range margins.  They do \emph{not} yet fix the accuracy
of bin comparisons.  Afterwards one may make a second,
arbitrarily thinner tile-boundary exclusion, refine the fixed
output bins, and increase the fixed integer \(j\).  These later
choices improve errors without changing the earlier coverage
requirements.

All these requirements can be met simultaneously with the
information estimate.  Initially average both over shifts and
over \(k\in[T/2,3T/4]\).  The first collar was chosen to have
expected loss small enough that its required coverage bound
holds with high probability.  Any later, finer collar can have
arbitrarily small expected loss.  For every fixed \(j\),
Lemma~\ref{lem:stage-information} gives an error tending to zero
in this joint average.  Markov's inequality therefore selects
common shifts and \(k=k_i\) satisfying the coverage bounds and
the desired comparison-error bounds.  The transverse grids in
\(q,u\) may be fixed using the grid convention of the preceding
section.  Every scale still satisfies
\eqref{eq:w-factor-scales}.

\paragraph{Geometry of the output comparisons.}
Write the coordinates of an assigned output chart as
\((q',w',u')\).  For fixed input \(q\), its patch image is obtained
by left multiplication in \(V\) from
\[
                     \Delta^{-m_0}\mathcal F
                            \ \times\ (\text{input }u\text{ interval}).
\]
Left multiplication by \(w_0u_0\) sends \(wu\) to
\begin{equation}\label{eq:separate-left-translation}
                   \bigl(w_0\Ad_{u_0}(w)\bigr)\,(u_0u).
\end{equation}
Thus the two plaque coordinates change separately.  The
parameters of this bounded left multiplication vary by
\(O(2^{-F}e^{CT})\) across the input \(q\)-cell.
The \(V\)-transverse diameter in \(q'\) is bounded by the same
quantity.  It is smaller than \(2^{-CS}\) for every fixed \(C\)
when \(i\) is large.
For a subpatch, the \(W\)-transverse coordinates are
\((q',u')\); their diameter is
\[
                         O(2^{-S}+2^{-F}e^{CT})
                                  =o(2^{-l_-}).
\]
Discard patches or subpatch comparisons crossing a coarse
output transverse wall.  Their sampling mass is \(o(1)\).

We apply Lemma~\ref{lem:patch-comparison} to \(V\) data on the
patches and \(W\) data on the subpatches.  The needed inner
products are as follows.  For a patch take collections of
\(w'\) and \(u'\) bins in product form; for each of its subpatches
use the same inner collection of \(w'\) bins.  These collections
are admitted over the entire corresponding transverse
projection, by \eqref{eq:separate-left-translation} and the
vanishing variation of its parameters.

We verify that these inner products lose as little sampling
mass as desired.  The extended charts describe one plaque sheet
throughout each of the ranges involved.  The transverse maps
are injective on a patch or subpatch: their images have a common
open plaque range, and injectivity follows from injectivity of
the original chart and the diagonal motion.  Vary a fine output
transverse coordinate across its cell and pull the path back.
Its input length is at most \(e^{CT}\) times the fine mesh side.
Because \(F,S\prec l_+\), this path stays in its input \(q\)- and
\(u\)-cells except at grid-wall collars of mass \(o(1)\).
Fixed input chart-boundary collars have arbitrarily small mass.
In tile coordinates the same pullback costs at most \(e^{C'T}\),
so it stays in the input tile away from any fixed thin collar
of that tile's boundary.  The second collar chosen above makes
this loss arbitrarily small.

Finally, away from these collars, the sample has a definite
margin in its admitted \(w'\) and \(u'\) ranges.  For \(w'\) the
margin depends on \(m_0\), the fixed charts, and the tile
collar thickness; for \(u'\) it follows from the fixed interval
margin.  Sufficiently fine fixed plaque bins are then wholly
admitted.  This proves the required plaque-bin trim estimate.
The collections can retain the supports of \(f,\phi\) about
the covered samples and a fixed \(U\) neighborhood about their
\(u'\) coordinates.  Thus the comparison errors can be made
arbitrarily small after the original coverage and range choices.

\paragraph{Comparison of the \(W\) tests.}
At a sampled output point, let \(P_V\) and \(P_W\) denote its
finite-\(i\) chart probabilities for \(V\) and \(W\).
Interpret \(f,\phi\) in the output \(w'\) coordinate using the
displacement formula \eqref{eq:product-displacement}.  Their
parameters may depend on the sample, but their norms and
continuity moduli are uniform on the finitely many chart ranges.
Let \(B_u\) denote the union of the inner \(u'\) bins for its
patch.  We claim that, on the covered samples,
\begin{equation}\label{eq:w-cross-product}
             P_V(f\,1_{B_u})P_W(\phi)
                    -P_V(\phi\,1_{B_u})P_W(f)
\end{equation}
has arbitrarily small expectation of its absolute value in the
limit \(i\to\infty\).

Here is the mass calculation.  Let \(A_p\) be the law of the
output plaque coordinates when sampling the input patch \(p\),
and let \(A_{p'}\) be the corresponding law for its subpatch.
The unnormalized form of Lemma~\ref{lem:patch-comparison},
first for \(V\) and then for \(W\), gives
\[
 \begin{aligned}
       c_p A_p(f)&\simeq P_V(f\,1_{B_u}),&
       c_p A_p(\phi)&\simeq P_V(\phi\,1_{B_u}),\\
       c_{p'} A_{p'}(f)&\simeq P_W(f),&
       c_{p'} A_{p'}(\phi)&\simeq P_W(\phi).
 \end{aligned}
\]
Let \(e_f(x),e_\phi(x),e'_f(x),e'_\phi(x)\) denote the absolute
errors in these four comparisons at the sampled point \(x\).
The factors \(c_p,c_{p'}\) are the probabilities of the admitted
plaque-bin ranges in the respective coarse output laws, hence
belong to \([0,1]\).  Each \(\simeq\) means an absolute error whose
sum weighted by the relevant patch masses is arbitrarily small.
The tests are supported in the common inner \(w'\) range, so
there is no additional \(w'\)-restriction in the right-hand
integrals.  Replacing coarse laws by the data laws uses the
plateau comparison.  These estimates remain true for the
sample-dependent test parameters by the uniform-test clause
of Lemma~\ref{lem:patch-comparison}.

At input, Lemma~\ref{lem:stage-information} compares the law of
the level-\((k+j)\) tile in \(p\) with its law in \(p'\).
At output these finer tiles have relative shapes
\(\Delta^{j-m_0}\mathcal F\).  Once \(m_0\) and the finite chart
ranges are fixed, taking the fixed integer \(j\) large makes
their diameters as small as prescribed.  The conversion to
\(w'\) coordinates has uniformly bounded distortion and is
independent of \(u\), by
\eqref{eq:separate-left-translation}.
Let \(\eta_i(x)\) be the \(\ell^1\) distance between the two
conditional level-\((k+j)\) tile laws for the patch and the
subpatch containing \(x\).  There is a uniform bound
\[
 \sigma_{i,j}
    \le C\,\diam(\Delta^{j-m_0}\mathcal F)+O(2^{-F}e^{CT})
\]
for the variation of \(w'\) within one such tile, including the
variation across its \(q\)-cell.  If \(h\) is either \(f\) or
\(\phi\), and \(\omega_h\) is a common continuity modulus for its
chart-coordinate versions, then
\begin{equation}\label{eq:tile-test-comparison}
       |A_p(h)-A_{p'}(h)|
                  \le \eta_i(x)+2\omega_h(\sigma_{i,j}).
\end{equation}
Indeed, replace \(h\) by one value on each finer tile.  The two
replacements cost at most \(\omega_h(\sigma_{i,j})\) each, and
the remaining finite-label difference is bounded by \(\eta_i(x)\).
This also proves the assertion for tests based at the sampled
point, since the bound is uniform over the test coefficients.
The expected value of \(\eta_i\) tends to zero for the selected
shifts and stages, by Lemma~\ref{lem:stage-information}.

Multiply the four mass comparisons and cancel by cross
multiplication, without dividing by \(c_p\) or \(c_{p'}\).
The result is \eqref{eq:w-cross-product}: the remaining expression
is bounded by
\[
 \begin{split}
 &e_f(x)+e_\phi(x)+e'_f(x)+e'_\phi(x)\\
 &\hspace{15mm}
       +2\eta_i(x)
       +2\omega_f(\sigma_{i,j})+2\omega_\phi(\sigma_{i,j}).
 \end{split}
\]
For this inequality, the term left after the four mass
comparisons is
\(c_pc_{p'}(A_p(f)A_{p'}(\phi)-A_p(\phi)A_{p'}(f))\);
apply \eqref{eq:tile-test-comparison} and use that all factors
are bounded by one.
Integrating over the covered samples bounds
the first four terms by the total trim and bin-comparison
errors, the fifth by the stage-information error, and the last
two by their uniform moduli.  Choose the trim accuracy and fixed
bins first, then choose the fixed \(j\) sufficiently large, and
finally let \(i\) tend to infinity.
This proves the claimed integrated bound for
\eqref{eq:w-cross-product}.
The number of patches may grow with \(i\); the absolute mass estimate
\eqref{eq:mass-error} charges each fine-cell/bin pair at most
once within a fixed chart comparison.

\paragraph{Removing the \(U\) restriction.}
By Lemma~\ref{lem:u-factor}, in the limiting \(V\) data the
\(w'\) and \(u'\) coordinates are independent.  Consequently,
for the finite-\(i\) probabilities,
\[
       P_V(f1_{B_u})-P_V(f)P_V(B_u)\longrightarrow0
\]
in probability on the charts under consideration, and the same
holds with \(\phi\).  To justify the varying \(B_u\), its bins
come from one fixed finite subdivision, and every bin boundary
has zero mass almost surely in the limiting data.  The asserted
convergence therefore holds simultaneously for all subsets of
that finite bin collection.  The displacement tests depend
continuously on the sample on the fixed inner chart ranges.
Output sampling has law \(\mu_i\), so the limiting independence
applies there even though the patches were chosen at input.
If several input charts are used, their finite multiplicity
only multiplies the error bound.

Thus \eqref{eq:w-cross-product} becomes
\begin{equation}\label{eq:w-cross-product-reduced}
 P_V(B_u)
       \bigl(P_V(f)P_W(\phi)-P_V(\phi)P_W(f)\bigr)
                            \simeq0.
\end{equation}

\paragraph{Denominators and passage to the projective limit.}
We now explain why \eqref{eq:w-cross-product-reduced} gives
\eqref{eq:w-ratio-target}.  On each of the fixed covered chart
ranges, \(B_u\) contains a fixed \(U\) neighborhood about the
sample.  Choose a continuous bump supported in that neighborhood
and positive at the sampled \(u'\) coordinate.  Its integral in
the limiting \(V\) probability is positive by identity support.
The same support property gives positivity of the limiting
integrals \(P_V(\phi)\) and \(P_W(\phi)\).
For any further prescribed exceptional probability there is,
therefore, a constant \(\kappa>0\) such that, for all large \(i\),
\begin{equation}\label{eq:w-denominators}
          P_V(B_u)\ge\kappa,\qquad
          P_V(\phi)\ge\kappa,\qquad
          P_W(\phi)\ge\kappa
\end{equation}
outside that exception on the covered samples.
This follows from convergence of the joint data in the finitely
many charts, using the fixed continuous bumps.

The order of choices is essential.  Fix the coverage, \(m_0\),
the output charts, and their test-range margins first.  Then
choose \(\kappa\) from these fixed bumps.  Only afterwards refine
the bin approximations and tile collars used for comparison,
and enlarge \(j\), to make the error in
\eqref{eq:w-cross-product-reduced} as small as required relative
to \(\kappa^3\).  The coverage collar and the fixed neighborhoods
giving the lower bounds need not change.  Division in
\eqref{eq:w-cross-product-reduced} is now legitimate and yields
arbitrarily small difference between
\[
                         \frac{P_V(f)}{P_V(\phi)}
                 \quad\hbox{and}\quad
                         \frac{P_W(f)}{P_W(\phi)}
\]
outside an arbitrarily small exceptional probability.
Explicitly, let \(\mathcal G_i\) be the covered set on which
\eqref{eq:w-denominators} holds, and denote the expression on the
left of \eqref{eq:w-cross-product-reduced} by \(D_i(x)\).
For every \(\delta>0\),
\[
 \mu_i\left(\mathcal G_i\cap
   \left\{\left|\frac{P_V(f)}{P_V(\phi)}
                -\frac{P_W(f)}{P_W(\phi)}\right|>\delta\right\}\right)
       \le \frac{1}{\delta\kappa^3}
                \int_{\mathcal G_i}|D_i(x)|\,\dd\mu_i(x).
\]
The right side can be made arbitrarily small with the previously
fixed coverage and denominators.

To see that this proves the assertion for the limiting
projective measures, suppose otherwise.  A countable choice of
tests and charts would then give a reference output chart and
a positive-probability set on which the two ratios differ by
more than some fixed positive number.  Restrict this set so
that the denominators have a positive lower bound and the
supports of the tests, together with a \(U\) neighborhood, have
inner margins in the chart.  Continuous displacement integrals
and the open-set inequality preserve a positive probability
of discrepancy in the sequence data.

Apply the preceding coverage construction with exceptional
probability much smaller than this discrepancy probability.
On samples covered by both the reference chart and an assigned
output chart, Proposition~\ref{prop:plaque-comparison} and
Lemma~\ref{lem:u-factor} identify their limiting ratios.
This also identifies their finite-\(i\) ratios up to an error
tending to zero outside a further small set: first impose
closed inner margins and lower bounds on the denominators,
then use joint convergence of the finitely many chart data.
The arbitrarily accurate comparison just proved in the
assigned charts contradicts persistence of the discrepancy
in the reference chart.

We have proved \eqref{eq:w-ratio-target}.  Taking countably many
tests and then exhausting compact \(W\) ranges gives
\(\nu^W_\omega\propto m^W_\omega\) almost surely.
Together with \eqref{eq:provisional-product}, this completes
the proof of Theorem~\ref{thm:product}.
